% Preamble for the G17 technical reference.
\usepackage{fontspec}
\usepackage{libertinus-otf}
\setmonofont{Inconsolatazi4}[Scale=MatchLowercase, ItalicFont=Inconsolatazi4, ItalicFeatures={FakeSlant=0.18}]
\usepackage{unicode-math}
\setmathfont{Libertinus Math}
\usepackage[final]{microtype}
\usepackage[english]{babel}
\usepackage[autostyle]{csquotes}
\MakeOuterQuote{"}

\usepackage[margin=1.05in, top=1in, bottom=1.1in]{geometry}
\usepackage{xcolor}
\definecolor{ink}{HTML}{1F2933}
\definecolor{accent}{HTML}{25569A}
\definecolor{apple}{HTML}{6B7280}
\definecolor{ours}{HTML}{1F5FA6}
\definecolor{native}{HTML}{B4462B}
\definecolor{rule}{HTML}{C9CED6}
\definecolor{codebg}{HTML}{F5F6F8}
\definecolor{hilite}{HTML}{F6D365}

\usepackage{enumitem}
\setlist{itemsep=2pt, topsep=0pt, parsep=0pt}
\setlist[itemize]{label=\textcolor{apple}{\textbullet}, leftmargin=1.4em}
\setlist[itemize,2]{label=\textcolor{apple}{\textendash}}

\usepackage{booktabs}
\usepackage{array}
\usepackage{longtable}
\usepackage{xltabular}
\renewcommand\tabularxcolumn[1]{>{\raggedright\arraybackslash}p{#1}}
\setlength{\LTpre}{9pt}
\setlength{\LTpost}{6pt}
\usepackage{etoolbox}
\newcommand{\tablesize}{\small}
% \widetable before a table with six or more columns sets it one size smaller.
\newcommand{\widetable}{\gdef\tablesize{\footnotesize}}
\AtBeginEnvironment{longtable}{\tablesize\renewcommand\arraystretch{1.12}\gdef\tablesize{\small}}
\AtBeginEnvironment{xltabular}{\tablesize\renewcommand\arraystretch{1.12}\gdef\tablesize{\small}}
\setlength{\extrarowheight}{1.5pt}
\setlength{\LTleft}{0pt}
\setlength{\LTright}{\fill}
\usepackage{seqsplit}

\usepackage{listings}
\lstset{basicstyle=\ttfamily\footnotesize, columns=fullflexible, keepspaces=true,
  breaklines=true, breakatwhitespace=false, backgroundcolor=\color{codebg},
  frame=none, xleftmargin=6pt, xrightmargin=6pt, framesep=4pt,
  aboveskip=8pt, belowskip=8pt, showstringspaces=false, upquote=true,
  extendedchars=true}
\lstdefinestyle{py}{language=Python, keywordstyle=\color{accent}\bfseries,
  commentstyle=\color{apple}\itshape, stringstyle=\color{native}}

\usepackage{tikz}
\usetikzlibrary{arrows.meta, positioning, fit, backgrounds, calc, shapes.misc, decorations.pathreplacing}
\usepackage{bytefield}
\bytefieldsetup{bitformatting=\footnotesize}
\usepackage{caption}
\captionsetup{font=small, labelfont=bf, format=plain, labelsep=period, justification=raggedright,
  singlelinecheck=false, margin=0pt}
\usepackage[most]{tcolorbox}

% Chapter-number sections in the opening (0.1, 0.2, ...).
\KOMAoptions{parskip=half-}
\setkomafont{disposition}{\normalfont\bfseries\color{ink}}
\setkomafont{part}{\normalfont\Huge\bfseries\color{ink}}
\setkomafont{chapter}{\normalfont\huge\bfseries\color{ink}}
\RedeclareSectionCommand[beforeskip=0pt, afterskip=1.2\baselineskip]{chapter}
\RedeclareSectionCommand[tocnumwidth=2.6em]{section}
\setkomafont{title}{\normalfont\bfseries\color{ink}}
\setkomafont{subtitle}{\normalfont\large\color{ink!80}}
\RedeclareSectionCommand[beforeskip=-1.4\baselineskip, afterskip=0.5\baselineskip]{section}
\RedeclareSectionCommand[beforeskip=-1.0\baselineskip, afterskip=0.3\baselineskip, font=\normalsize\bfseries]{subsection}
\setcounter{secnumdepth}{1}
\setcounter{tocdepth}{1}
\counterwithout{figure}{chapter}
\renewcommand{\thefigure}{\thechapter.\arabic{figure}}
\counterwithin{figure}{chapter}
\counterwithin{table}{chapter}
\captionsetup[table]{position=top, skip=4pt}
% A table's caption row: \caption{...}\label{tab:...}\\ before \toprule. The header is repeated under a
% "continued" line when a long table breaks across pages.
\newcommand{\tablecontinued}[1]{\multicolumn{#1}{@{}l@{}}{\footnotesize\itshape\tablename~\thetable, continued}\\}

\usepackage[hyphens]{url}
\usepackage[colorlinks=true, linkcolor=accent, urlcolor=accent, citecolor=accent, linktoc=page,
  pdftitle={The G17 technical reference}, bookmarksnumbered=true]{hyperref}
\usepackage{bookmark}
% Cross-references: \cref{x} prints "section 10.4", "chapter 9", "appendix C" or "figure 3.1" from the label's
% own type, so the word cannot go stale; \Cref at the start of a sentence; \crefrange{a}{b} for "sections 2.4 to 2.6".
\usepackage[nameinlink, noabbrev]{cleveref}
\crefname{part}{Part}{Parts}
\crefname{chapter}{chapter}{chapters}
\crefname{section}{section}{sections}
\crefname{appendix}{appendix}{appendices}
\crefname{figure}{figure}{figures}
\crefname{table}{table}{tables}
\newcommand{\crefrangeconjunction}{ to~}
\pdfstringdefDisableCommands{\def\_{_}\def\code#1{#1}\def\op#1{op#1}\def\Op#1{op#1}\def\opn#1{op#1}\def\val#1{#1}}

% --- Semantic macros -------------------------------------------------------
% Inline code may break after / . _ , ; : ( and before =, so long paths and
% identifiers do not run into the margin.
\usepackage{luacode}
\begin{luacode*}
function g17_code(s)
  local long = #s > 16
  s = s:gsub('"', '\0Q')
  s = s:gsub('\\_ ?', '\0U'):gsub('\\%& ?', '\0A'):gsub('\\%# ?', '\0H'):gsub('\\%$ ?', '\0D')
  s = s:gsub('\\%% ?', '\0P'):gsub('\\{ ?', '\0L'):gsub('\\} ?', '\0R')
  s = s:gsub('\\textbackslash ?{}', '\0B'):gsub('\\textasciicircum ?{}', '\0C'):gsub('\\textasciitilde ?{}', '\0T')
  s = s:gsub('([/,;])', '%1\0Z')
  if long then
    s = s:gsub('%z(U)', '\0V')
    s = s:gsub('([%.|(])', '%1\0Z')
    s = s:gsub('%zZ(%a%a?%a?%a?)$', '%1')   -- no break before a file extension
  end
  if #s > 30 then s = s:gsub('(%l%l%l)(%u%l%l)', '%1\0Z%2') end
  local m = {U='\\textunderscore{}', V='\\textunderscore\\allowbreak{}', Q='\\textquotedbl{}', A='\\&', H='\\#', D='\\$', P='\\%', L='\\{', R='\\}',
             B='\\textbackslash{}', C='\\textasciicircum{}', T='\\textasciitilde{}', Z='\\allowbreak{}'}
  s = s:gsub('%z(%u)', m)
  tex.sprint(s)
end
\end{luacode*}
\DeclareRobustCommand{\code}[1]{\texttt{\exhyphenpenalty=10000\hyphenchar\font=-1 \directlua{g17_code("\luaescapestring{\detokenize{#1}}")}}}
\emergencystretch=3em
\widowpenalty=10000
\clubpenalty=10000
\usepackage{needspace}
% Keep a heading with at least a few lines of what follows it.
\AddToHook{cmd/section/before}{\Needspace{7\baselineskip}}
\AddToHook{cmd/subsection/before}{\Needspace{5\baselineskip}}
% A bold run-in line that introduces a list or table stays with it.
\newcommand{\leadhead}[1]{\Needspace{4\baselineskip}\noindent\textbf{#1}\par\nopagebreak}
\makeatletter
\setlength{\@fptop}{0pt}
\makeatother
\renewcommand{\floatpagefraction}{0.75}
\renewcommand{\topfraction}{0.85}
\renewcommand{\textfraction}{0.1}
\newcommand{\hash}[1]{\texttt{\seqsplit{#1}}}

% Opcodes name themselves. \op{5107} prints the glossary name and the number, linked to the opcode's row in
% appendix C; \Op{5107} capitalises the name at a sentence start; \opn{5107} prints the linked number alone (a
% table cell beside a name column, or a deliberately shortened name); \opdef{5107} makes the row. The names come
% from opnames.tex, generated from isa/g17-opcode-glossary.json by tools/g17docs.py --write-opnames. An opcode
% with no glossary entry prints its number only, and its paragraph says "no glossary entry".
\newcommand{\opnamedef}[3]{\expandafter\def\csname opname@#1\endcsname{#2}%
  \expandafter\def\csname Opname@#1\endcsname{#3}}
\DeclareRobustCommand{\opn}[1]{\hyperlink{op:#1}{\textcolor{ink}{\texttt{op#1}}}}
\DeclareRobustCommand{\op}[1]{\ifcsname opname@#1\endcsname\csname opname@#1\endcsname\ (\opn{#1})\else\opn{#1}\fi}
\DeclareRobustCommand{\Op}[1]{\ifcsname Opname@#1\endcsname\csname Opname@#1\endcsname\ (\opn{#1})\else\opn{#1}\fi}
\newcommand{\opdef}[1]{\hypertarget{op:#1}{\texttt{op#1}}}
% GENERATED by tools/g17docs.py --write-opnames from isa/g17-opcode-glossary.json; do not edit.
% \op{N} prints the name and the linked number; an opcode with no entry prints the number only.
\opnamedef{0}{uniform register file, decoder tag}{Uniform register file, decoder tag}
\opnamedef{4}{message staging file, decoder tag}{Message staging file, decoder tag}
\opnamedef{423}{32-bit AND with immediate}{32-bit AND with immediate}
\opnamedef{424}{32-bit AND}{32-bit AND}
\opnamedef{426}{16-bit AND with immediate, 32-bit result}{16-bit AND with immediate, 32-bit result}
\opnamedef{428}{16-bit AND, 32-bit result}{16-bit AND, 32-bit result}
\opnamedef{435}{16-bit AND with immediate}{16-bit AND with immediate}
\opnamedef{437}{16-bit AND}{16-bit AND}
\opnamedef{441}{32-bit AND with immediate into the uniform/staging file}{32-bit AND with immediate into the uniform/staging file}
\opnamedef{447}{threadgroup barrier}{Threadgroup barrier}
\opnamedef{450}{call}{Call}
\opnamedef{458}{conditional back-edge branch}{Conditional back-edge branch}
\opnamedef{462}{skip branch, no active lanes}{Skip branch, no active lanes}
\opnamedef{465}{population count}{Population count}
\opnamedef{554}{move immediate, two-operand form}{Move immediate, two-operand form}
\opnamedef{555}{16-bit move immediate into a half}{16-bit move immediate into a half}
\opnamedef{577}{exec pop, restore mask}{Exec pop, restore mask}
\opnamedef{579}{while-exec, set mask in place}{While-exec, set mask in place}
\opnamedef{580}{return}{Return}
\opnamedef{582}{exec push, if, from flag}{Exec push, if, from flag}
\opnamedef{586}{32-bit register move}{32-bit register move}
\opnamedef{590}{16-bit register move}{16-bit register move}
\opnamedef{592}{uniform / staging write}{Uniform / staging write}
\opnamedef{593}{publish sibling, function unknown}{Publish sibling, function unknown}
\opnamedef{595}{fragment spill store}{Fragment spill store}
\opnamedef{612}{range mask, immediate width}{Range mask, immediate width}
\opnamedef{615}{range mask, register width}{Range mask, register width}
\opnamedef{621}{range mask, 16-bit source}{Range mask, 16-bit source}
\opnamedef{684}{end of program}{End of program}
\opnamedef{776}{half add with float immediate}{Half add with float immediate}
\opnamedef{798}{fp16 fma, register form}{Fp16 fma, register form}
\opnamedef{838}{legacy 8×8 simdgroup MMA, f16}{Legacy 8×8 simdgroup MMA, f16}
\opnamedef{839}{legacy 8×8 simdgroup multiply, f16}{Legacy 8×8 simdgroup multiply, f16}
\opnamedef{904}{fp32 saturate}{Fp32 saturate}
\opnamedef{998}{fp32 add}{Fp32 add}
\opnamedef{1001}{fp32 add, rare alternate encoding}{Fp32 add, rare alternate encoding}
\opnamedef{1004}{half to fp32 widening}{Half to fp32 widening}
\opnamedef{1006}{fp32 add with float immediate}{Fp32 add with float immediate}
\opnamedef{1014}{fp32 add, half result}{Fp32 add, half result}
\opnamedef{1015}{fp32 plus half, half result}{Fp32 plus half, half result}
\opnamedef{1016}{fp32 to half narrowing}{Fp32 to half narrowing}
\opnamedef{1038}{fp32 add with immediate into the uniform/staging file}{Fp32 add with immediate into the uniform/staging file}
\opnamedef{1048}{fp32 to bfloat16 narrowing, unconfirmed}{Fp32 to bfloat16 narrowing, unconfirmed}
\opnamedef{1061}{fp32 add, common alternate encoding}{Fp32 add, common alternate encoding}
\opnamedef{1062}{fine x-derivative, saturated}{Fine x-derivative, saturated}
\opnamedef{1272}{hardware exp2}{Hardware exp2}
\opnamedef{2190}{fp32 fused multiply-add}{Fp32 fused multiply-add}
\opnamedef{2192}{fp32 FMA, literal addend}{Fp32 FMA, literal addend}
\opnamedef{2198}{fp32 FMA, literal multiplicand}{Fp32 FMA, literal multiplicand}
\opnamedef{2200}{fp32 FMA with two float immediates}{Fp32 FMA with two float immediates}
\opnamedef{2238}{fp32 fused multiply-add, alternate encoding}{Fp32 fused multiply-add, alternate encoding}
\opnamedef{2240}{FMA-family variant, narrow operand, unconfirmed}{FMA-family variant, narrow operand, unconfirmed}
\opnamedef{2241}{fused multiply-add, schedclass-87 base form}{Fused multiply-add, schedclass-87 base form}
\opnamedef{2318}{fp32 FMA into the uniform/staging file}{Fp32 FMA into the uniform/staging file}
\opnamedef{2433}{fused multiply-add, narrowed form}{Fused multiply-add, narrowed form}
\opnamedef{2570}{hardware log2}{Hardware log2}
\opnamedef{2842}{legacy 8×8 simdgroup MMA, fp32}{Legacy 8×8 simdgroup MMA, fp32}
\opnamedef{2844}{legacy 8×8 simdgroup multiply, fp32}{Legacy 8×8 simdgroup multiply, fp32}
\opnamedef{2862}{legacy 8×8 simdgroup MMA, f16 in fp32 out}{Legacy 8×8 simdgroup MMA, f16 in fp32 out}
\opnamedef{2902}{legacy 8×8 simdgroup MMA, bf16}{Legacy 8×8 simdgroup MMA, bf16}
\opnamedef{3290}{fp32 multiply}{Fp32 multiply}
\opnamedef{3291}{fp32 times half multiply}{Fp32 times half multiply}
\opnamedef{3298}{fp32 op with float immediate, unconfirmed}{Fp32 op with float immediate, unconfirmed}
\opnamedef{3306}{fp32 multiply, half result}{Fp32 multiply, half result}
\opnamedef{3322}{fp32 multiply into the uniform/staging file}{Fp32 multiply into the uniform/staging file}
\opnamedef{3653}{16-bit reciprocal into the uniform/staging file}{16-bit reciprocal into the uniform/staging file}
\opnamedef{3658}{hardware reciprocal}{Hardware reciprocal}
\opnamedef{3770}{round to nearest even}{Round to nearest even}
\opnamedef{3786}{floor}{Floor}
\opnamedef{3818}{fp32 truncate toward zero}{Fp32 truncate toward zero}
\opnamedef{3850}{rsqrt seed}{Rsqrt seed}
\opnamedef{3946}{trig-family op, function unmeasured}{Trig-family op, function unmeasured}
\opnamedef{3978}{sqrt-safe rsqrt}{Sqrt-safe rsqrt}
\opnamedef{5098}{tensor MMA, fp32 × fp32, accumulate}{Tensor MMA, fp32 × fp32, accumulate}
\opnamedef{5099}{tensor MMA, fp32 × fp32, no accumulator}{Tensor MMA, fp32 × fp32, no accumulator}
\opnamedef{5100}{tensor MMA, fp32 A × 16-bit B, accumulate}{Tensor MMA, fp32 A × 16-bit B, accumulate}
\opnamedef{5101}{tensor MMA, fp32 A × 16-bit B, no accumulator}{Tensor MMA, fp32 A × 16-bit B, no accumulator}
\opnamedef{5104}{tensor MMA, 16-bit A × fp32 B, accumulate}{Tensor MMA, 16-bit A × fp32 B, accumulate}
\opnamedef{5105}{tensor MMA, 16-bit A × fp32 B, no accumulator}{Tensor MMA, 16-bit A × fp32 B, no accumulator}
\opnamedef{5106}{tensor MMA, 16-bit × 16-bit, accumulate}{Tensor MMA, 16-bit × 16-bit, accumulate}
\opnamedef{5107}{tensor MMA, 16-bit × 16-bit, no accumulator}{Tensor MMA, 16-bit × 16-bit, no accumulator}
\opnamedef{5108}{table-only narrow-accumulator MMA, first}{Table-only narrow-accumulator MMA, first}
\opnamedef{5143}{table-only narrow-accumulator MMA, last, no accumulator}{Table-only narrow-accumulator MMA, last, no accumulator}
\opnamedef{9320}{fp32 to 32-bit integer}{Fp32 to 32-bit integer}
\opnamedef{9321}{half to 32-bit integer}{Half to 32-bit integer}
\opnamedef{9325}{half to 16-bit integer}{Half to 16-bit integer}
\opnamedef{9328}{fp32 to u32 into the uniform/staging file}{Fp32 to u32 into the uniform/staging file}
\opnamedef{9329}{half to 32-bit integer into the uniform/staging file}{Half to 32-bit integer into the uniform/staging file}
\opnamedef{9683}{half compare, chained into a flag}{Half compare, chained into a flag}
\opnamedef{9700}{fp32 compare-select, fmax/fmin}{Fp32 compare-select, fmax/fmin}
\opnamedef{9721}{32-bit compare-select, immediate arms}{32-bit compare-select, immediate arms}
\opnamedef{9748}{float compare-select, fselect form}{Float compare-select, fselect form}
\opnamedef{9749}{float compare-select, csel form}{Float compare-select, csel form}
\opnamedef{9832}{half compare-select}{Half compare-select}
\opnamedef{9842}{half compare-select against minifloat}{Half compare-select against minifloat}
\opnamedef{9986}{find most significant bit}{Find most significant bit}
\opnamedef{9989}{msb bit scan, 16-bit result}{Msb bit scan, 16-bit result}
\opnamedef{9996}{64-bit no-return atomic, max/min}{64-bit no-return atomic, max/min}
\opnamedef{10019}{per-lane-address no-return atomic}{Per-lane-address no-return atomic}
\opnamedef{10022}{uniform-address no-return atomic, 32-bit}{Uniform-address no-return atomic, 32-bit}
\opnamedef{10023}{uniform-address no-return compare-exchange}{Uniform-address no-return compare-exchange}
\opnamedef{10090}{per-lane-address atomic, returning}{Per-lane-address atomic, returning}
\opnamedef{10091}{per-lane-address compare-exchange, returning}{Per-lane-address compare-exchange, returning}
\opnamedef{10094}{uniform-address atomic, returning old}{Uniform-address atomic, returning old}
\opnamedef{10095}{uniform-address compare-exchange, returning}{Uniform-address compare-exchange, returning}
\opnamedef{10239}{32-bit saturating add}{32-bit saturating add}
\opnamedef{10267}{64-bit add}{64-bit add}
\opnamedef{10279}{32-bit add immediate}{32-bit add immediate}
\opnamedef{10282}{32-bit add}{32-bit add}
\opnamedef{10283}{add, 32-bit plus zero-extended 16-bit}{Add, 32-bit plus zero-extended 16-bit}
\opnamedef{10285}{add 16-bit to 32-bit}{Add 16-bit to 32-bit}
\opnamedef{10286}{16-bit plus 16-bit, 32-bit result}{16-bit plus 16-bit, 32-bit result}
\opnamedef{10289}{16-bit add immediate}{16-bit add immediate}
\opnamedef{10298}{32-bit add immediate into the uniform/staging file}{32-bit add immediate into the uniform/staging file}
\opnamedef{10306}{32-bit add into the uniform/staging file}{32-bit add into the uniform/staging file}
\opnamedef{10369}{32-bit compare with immediate into a flag}{32-bit compare with immediate into a flag}
\opnamedef{10372}{16-bit compare into a flag}{16-bit compare into a flag}
\opnamedef{10374}{16-bit register compare into a flag}{16-bit register compare into a flag}
\opnamedef{10378}{32-bit compare with immediate into a flag, short form}{32-bit compare with immediate into a flag, short form}
\opnamedef{10384}{tensor MMA, int8 widening to int32}{Tensor MMA, int8 widening to int32}
\opnamedef{10385}{tensor MMA, int8 widening, no accumulator}{Tensor MMA, int8 widening, no accumulator}
\opnamedef{10789}{32 × immediate to 64-bit product}{32 × immediate to 64-bit product}
\opnamedef{10790}{64-bit address multiply-add}{64-bit address multiply-add}
\opnamedef{10793}{32 × 32 to 64-bit product}{32 × 32 to 64-bit product}
\opnamedef{10795}{32 × 32 + 32 to 64-bit multiply-add}{32 × 32 + 32 to 64-bit multiply-add}
\opnamedef{10823}{multiply by immediate, 32-bit addend}{Multiply by immediate, 32-bit addend}
\opnamedef{10824}{multiply by immediate, 16-bit addend}{Multiply by immediate, 16-bit addend}
\opnamedef{10825}{32-bit multiply}{32-bit multiply}
\opnamedef{10826}{32-bit multiply-add}{32-bit multiply-add}
\opnamedef{10829}{multiply-add, 16-bit multiplier}{Multiply-add, 16-bit multiplier}
\opnamedef{10830}{multiply-add, 16-bit multiplier and addend}{Multiply-add, 16-bit multiplier and addend}
\opnamedef{10831}{16-bit multiply by immediate plus immediate}{16-bit multiply by immediate plus immediate}
\opnamedef{10835}{multiply-add, 16-bit multiplicand}{Multiply-add, 16-bit multiplicand}
\opnamedef{10836}{multiply-add, 16-bit multiplicand and addend}{Multiply-add, 16-bit multiplicand and addend}
\opnamedef{10838}{signed-by-unsigned 16-bit multiply-add}{Signed-by-unsigned 16-bit multiply-add}
\opnamedef{10858}{16-bit multiply by immediate plus immediate, 16-bit result}{16-bit multiply by immediate plus immediate, 16-bit result}
\opnamedef{10860}{16-bit multiply by immediate plus register}{16-bit multiply by immediate plus register}
\opnamedef{10879}{multiply by immediate plus immediate into the uniform/staging file}{Multiply by immediate plus immediate into the uniform/staging file}
\opnamedef{10883}{32-bit multiply into the uniform/staging file}{32-bit multiply into the uniform/staging file}
\opnamedef{10904}{MMA bit-flip mutation artifact}{MMA bit-flip mutation artifact}
\opnamedef{10909}{texture write}{Texture write}
\opnamedef{11059}{multiply by immediate minus immediate}{Multiply by immediate minus immediate}
\opnamedef{11060}{multiply by immediate minus register}{Multiply by immediate minus register}
\opnamedef{11063}{multiply-subtract}{Multiply-subtract}
\opnamedef{11120}{multiply-add into the uniform/staging file}{Multiply-add into the uniform/staging file}
\opnamedef{11179}{32-bit integer to fp32}{32-bit integer to fp32}
\opnamedef{11180}{16-bit integer to fp32}{16-bit integer to fp32}
\opnamedef{11182}{u32 to half}{U32 to half}
\opnamedef{11183}{16-bit integer to half}{16-bit integer to half}
\opnamedef{11185}{integer to float into the uniform/staging file}{Integer to float into the uniform/staging file}
\opnamedef{11190}{32-bit NOT}{32-bit NOT}
\opnamedef{11193}{NOT to 16-bit result}{NOT to 16-bit result}
\opnamedef{11196}{32-bit NOT into the uniform/staging file}{32-bit NOT into the uniform/staging file}
\opnamedef{11310}{32-bit compare, AND-chained into a flag}{32-bit compare, AND-chained into a flag}
\opnamedef{11313}{32-bit compare, OR-chained into a flag}{32-bit compare, OR-chained into a flag}
\opnamedef{11314}{32-bit compare, AND-chained into a flag with false arm}{32-bit compare, AND-chained into a flag with false arm}
\opnamedef{11322}{16-bit compare, chained into a flag}{16-bit compare, chained into a flag}
\opnamedef{11362}{32-bit compare-select against immediate}{32-bit compare-select against immediate}
\opnamedef{11364}{integer clamp}{Integer clamp}
\opnamedef{11365}{32-bit compare-select, three registers}{32-bit compare-select, three registers}
\opnamedef{11372}{32-bit register compare-select, immediate arms}{32-bit register compare-select, immediate arms}
\opnamedef{11375}{32-bit compare-select, four registers}{32-bit compare-select, four registers}
\opnamedef{11392}{16-bit compare-select against immediate, 32-bit result}{16-bit compare-select against immediate, 32-bit result}
\opnamedef{11415}{16-bit compare float select}{16-bit compare float select}
\opnamedef{11456}{32-bit compare against immediate, select to 16-bit}{32-bit compare against immediate, select to 16-bit}
\opnamedef{11458}{32-bit compare against immediate, select to 16-bit, minifloat false arm}{32-bit compare against immediate, select to 16-bit, minifloat false arm}
\opnamedef{11460}{32-bit compare against immediate, select to 16-bit, register false arm}{32-bit compare against immediate, select to 16-bit, register false arm}
\opnamedef{11461}{32-bit compare against immediate to half boolean}{32-bit compare against immediate to half boolean}
\opnamedef{11462}{32-bit register compare to 16-bit result}{32-bit register compare to 16-bit result}
\opnamedef{11466}{32-bit register compare, select to 16-bit}{32-bit register compare, select to 16-bit}
\opnamedef{11471}{32-bit register compare to half boolean}{32-bit register compare to half boolean}
\opnamedef{11472}{32-bit vs 16-bit compare to 16-bit result}{32-bit vs 16-bit compare to 16-bit result}
\opnamedef{11482}{16-bit compare-select against immediate}{16-bit compare-select against immediate}
\opnamedef{11491}{16-bit compare against immediate to half boolean}{16-bit compare against immediate to half boolean}
\opnamedef{11507}{16-bit compare-select, four registers}{16-bit compare-select, four registers}
\opnamedef{11564}{texture-coordinate staging write}{Texture-coordinate staging write}
\opnamedef{11624}{32-bit saturating subtract}{32-bit saturating subtract}
\opnamedef{11655}{64-bit subtract immediate}{64-bit subtract immediate}
\opnamedef{11666}{32-bit subtract immediate}{32-bit subtract immediate}
\opnamedef{11667}{32-bit subtract}{32-bit subtract}
\opnamedef{11689}{32-bit subtract immediate into the uniform/staging file}{32-bit subtract immediate into the uniform/staging file}
\opnamedef{11701}{threadgroup atomic, no return}{Threadgroup atomic, no return}
\opnamedef{11765}{threadgroup atomic, returning old}{Threadgroup atomic, returning old}
\opnamedef{11842}{move immediate, three-operand form}{Move immediate, three-operand form}
\opnamedef{11994}{ray-tracing block load, 32-bit}{Ray-tracing block load, 32-bit}
\opnamedef{12073}{device load, register address}{Device load, register address}
\opnamedef{12074}{device load, 16-bit register address}{Device load, 16-bit register address}
\opnamedef{12151}{imageblock load, 32-bit}{Imageblock load, 32-bit}
\opnamedef{12352}{MMA operand load, space unconfirmed}{MMA operand load, space unconfirmed}
\opnamedef{12364}{threadgroup-memory load}{Threadgroup-memory load}
\opnamedef{12397}{load, a further form}{Load, a further form}
\opnamedef{12646}{16-bit load}{16-bit load}
\opnamedef{12655}{packed half2 load}{Packed half2 load}
\opnamedef{12656}{int8 tensor load, one word}{Int8 tensor load, one word}
\opnamedef{12657}{masked int8 tensor load}{Masked int8 tensor load}
\opnamedef{12673}{plain 64-bit load, pair address}{Plain 64-bit load, pair address}
\opnamedef{12674}{tensor load, half/bf16}{Tensor load, half/bf16}
\opnamedef{12675}{masked tensor load, half/bf16}{Masked tensor load, half/bf16}
\opnamedef{12682}{plain 32-bit load}{Plain 32-bit load}
\opnamedef{12688}{immediate-address one-word load}{Immediate-address one-word load}
\opnamedef{12691}{plain 64-bit load}{Plain 64-bit load}
\opnamedef{12692}{tensor load, fp32}{Tensor load, fp32}
\opnamedef{12693}{masked tensor load, fp32}{Masked tensor load, fp32}
\opnamedef{12697}{immediate-address two-word load}{Immediate-address two-word load}
\opnamedef{12700}{three-component vector load}{Three-component vector load}
\opnamedef{12706}{immediate-address three-word load}{Immediate-address three-word load}
\opnamedef{12709}{128-bit vector load}{128-bit vector load}
\opnamedef{12715}{immediate-address four-word load}{Immediate-address four-word load}
\opnamedef{13075}{imageblock store, 32-bit}{Imageblock store, 32-bit}
\opnamedef{13276}{dequantized-block store, likely threadgroup}{Dequantized-block store, likely threadgroup}
\opnamedef{13288}{threadgroup-memory store}{Threadgroup-memory store}
\opnamedef{13483}{two-byte filler}{Two-byte filler}
\opnamedef{13574}{32-bit OR with immediate, seven-bit destination}{32-bit OR with immediate, seven-bit destination}
\opnamedef{13575}{32-bit OR}{32-bit OR}
\opnamedef{13576}{OR with 16-bit operand}{OR with 16-bit operand}
\opnamedef{13588}{16-bit OR}{16-bit OR}
\opnamedef{13592}{32-bit OR with immediate into the uniform/staging file}{32-bit OR with immediate into the uniform/staging file}
\opnamedef{13597}{16-bit OR into the uniform/staging file}{16-bit OR into the uniform/staging file}
\opnamedef{13618}{fp8 pack, vector, from fp32/half/bf16}{Fp8 pack, vector, from fp32/half/bf16}
\opnamedef{13620}{fp8 pack, scalar, from fp32}{Fp8 pack, scalar, from fp32}
\opnamedef{13754}{camera/video kernel op, function unknown}{Camera/video kernel op, function unknown}
\opnamedef{13856}{quad prefix sum, fp32}{Quad prefix sum, fp32}
\opnamedef{13944}{quad shuffle-xor, immediate mask}{Quad shuffle-xor, immediate mask}
\opnamedef{13945}{quad shuffle-xor, runtime mask}{Quad shuffle-xor, runtime mask}
\opnamedef{14047}{bit reverse}{Bit reverse}
\opnamedef{14059}{special-register read}{Special-register read}
\opnamedef{14060}{special-register read, 16-bit}{Special-register read, 16-bit}
\opnamedef{14061}{argument read, binding-table word}{Argument read, binding-table word}
\opnamedef{14099}{flag to 16-bit register}{Flag to 16-bit register}
\opnamedef{14147}{bit interleave, Morton encode}{Bit interleave, Morton encode}
\opnamedef{14156}{device fence}{Device fence}
\opnamedef{14157}{SIMD shuffle, constant lane, broadcast}{SIMD shuffle, constant lane, broadcast}
\opnamedef{14158}{SIMD shuffle, register lane}{SIMD shuffle, register lane}
\opnamedef{14169}{SIMD shuffle-xor, immediate mask}{SIMD shuffle-xor, immediate mask}
\opnamedef{14170}{SIMD shuffle-xor, runtime mask}{SIMD shuffle-xor, runtime mask}
\opnamedef{14208}{simdgroup active-thread mask}{Simdgroup active-thread mask}
\opnamedef{14211}{SIMD boolean reduce, base form}{SIMD boolean reduce, base form}
\opnamedef{14214}{SIMD boolean reduce}{SIMD boolean reduce}
\opnamedef{14391}{32-bit shift left immediate}{32-bit shift left immediate}
\opnamedef{14392}{32-bit shift left by register}{32-bit shift left by register}
\opnamedef{14394}{16-bit shift left immediate, 32-bit result}{16-bit shift left immediate, 32-bit result}
\opnamedef{14423}{shift left, a further form}{Shift left, a further form}
\opnamedef{14661}{texture sample}{Texture sample}
\opnamedef{15813}{texture read}{Texture read}
\opnamedef{16805}{32-bit arithmetic shift right immediate}{32-bit arithmetic shift right immediate}
\opnamedef{16842}{SIMD sum, fp32}{SIMD sum, fp32}
\opnamedef{16858}{SIMD max, fp32}{SIMD max, fp32}
\opnamedef{17013}{32-bit shift right immediate}{32-bit shift right immediate}
\opnamedef{17014}{32-bit shift right by register}{32-bit shift right by register}
\opnamedef{17016}{16-bit shift right, 32-bit result}{16-bit shift right, 32-bit result}
\opnamedef{17043}{16-bit shift right}{16-bit shift right}
\opnamedef{17193}{16-bit store}{16-bit store}
\opnamedef{17199}{16-bit store at an immediate address}{16-bit store at an immediate address}
\opnamedef{17202}{word store at word index}{Word store at word index}
\opnamedef{17229}{plain 32-bit store}{Plain 32-bit store}
\opnamedef{17235}{store, sub-form 01, texel/word-slot}{Store, sub-form 01, texel/word-slot}
\opnamedef{17244}{two-component store}{Two-component store}
\opnamedef{17253}{immediate-address three-word store}{Immediate-address three-word store}
\opnamedef{17256}{128-bit vector store}{128-bit vector store}
\opnamedef{17257}{tensor store}{Tensor store}
\opnamedef{17258}{masked tensor store}{Masked tensor store}
\opnamedef{17262}{immediate-address four-word store}{Immediate-address four-word store}
\opnamedef{17631}{convert-family base, class 11, unused}{Convert-family base, class 11, unused}
\opnamedef{17632}{fp8 unpack, scalar, to fp32}{Fp8 unpack, scalar, to fp32}
\opnamedef{17633}{convert-family base, class 34, unused}{Convert-family base, class 34, unused}
\opnamedef{17634}{fp8 unpack, vector, to fp32 or half}{Fp8 unpack, vector, to fp32 or half}
\opnamedef{17635}{convert-family base, store-shaped A, unused}{Convert-family base, store-shaped A, unused}
\opnamedef{17636}{convert-family, store-shaped A, unused}{Convert-family, store-shaped A, unused}
\opnamedef{17637}{convert-family base, store-shaped B, unused}{Convert-family base, store-shaped B, unused}
\opnamedef{17638}{convert-family, store-shaped B, unused}{Convert-family, store-shaped B, unused}
\opnamedef{17639}{convert-family base, class 0, unused}{Convert-family base, class 0, unused}
\opnamedef{17640}{convert-family, class 0 value-producing, unused}{Convert-family, class 0 value-producing, unused}
\opnamedef{17641}{convert-family base, class 12, unused}{Convert-family base, class 12, unused}
\opnamedef{17642}{fp8 unpack to bf16}{Fp8 unpack to bf16}
\opnamedef{17771}{32-bit XOR}{32-bit XOR}
\opnamedef{17772}{XOR with 16-bit second operand}{XOR with 16-bit second operand}
\opnamedef{17774}{XOR with 16-bit first operand}{XOR with 16-bit first operand}
\opnamedef{17775}{16-bit XOR, 32-bit result}{16-bit XOR, 32-bit result}
\opnamedef{17784}{16-bit XOR}{16-bit XOR}
\opnamedef{17788}{32-bit XOR with immediate into the uniform/staging file}{32-bit XOR with immediate into the uniform/staging file}
\opnamedef{17789}{scalar spill store, 32-bit XOR into the staging file}{Scalar spill store, 32-bit XOR into the staging file}
\opnamedef{17793}{16-bit XOR into the uniform/staging file}{16-bit XOR into the uniform/staging file}

% Numbers quoted from generated artifacts: \val{d1} prints the value from values.tex, generated by
% tools/g17docs.py --write-values (its VALUES list names each artifact). An unknown key prints ?? and warns.
\newcommand{\valdef}[2]{\expandafter\def\csname val@#1\endcsname{#2}}
\DeclareRobustCommand{\val}[1]{\ifcsname val@#1\endcsname\csname val@#1\endcsname\else
  \textbf{??}\PackageWarning{g17}{Undefined value #1}\fi}
% GENERATED by tools/g17docs.py --write-values from the artifacts named in its VALUES; do not edit.
\valdef{d1}{7,546}% D1 structural forms
\valdef{d2}{1,118}% D2 encoded forms
\valdef{d3-any}{340}% D3 forms, any own instrument
\valdef{d3-verified}{89}% D3 verified class
\valdef{dispatched-unverified}{472}% dispatched forms outside the verified class
\valdef{dispatched-other-evidence}{236}% of those, with other semantic evidence
\valdef{dispatched-no-evidence}{236}% dispatched forms with no semantic evidence
\valdef{admitted-opcodes}{6,719}% decoder-admitted opcode records
\valdef{named-instructions}{280}% glossary opcodes that are instructions (decoder tags excluded)
\valdef{executed-instructions}{193}% glossary opcodes confirmed by execution
\valdef{renumber-new-opcodes}{17,905}% 26A434 instruction table entries
\valdef{renumber-old-registers}{3,567}% 25G83 register table entries
\valdef{renumber-new-registers}{3,685}% 26A434 register table entries
\valdef{renumber-corpus-instructions}{184,349}% corpus instructions re-decoded
\valdef{renumber-witnessed}{3,244}% opcodes anchored by decoded encodings
\valdef{renumber-aligned}{11,583}% opcodes mapped by descriptor alignment
\valdef{renumber-ambiguous}{934}% opcodes left unmapped as ambiguous
\valdef{renumber-tags}{114}% fitted immediate-tag entries


% Machine-model citations: \mm{25.91} prints the section number and links to it.
\newcommand{\mmbase}{https://github.com/sbryngelson/AGXForge/blob/main/docs/g17-tensorops-machine-model.md}
\makeatletter
\newcommand{\mmdef}[2]{\expandafter\def\csname mm@#1\endcsname{#2}}
\def\mm@show#1/#2\@nil{#1}
% mm-anchors.tex is generated from the machine model's headings on every build (tools/g17docs.py
% --write-mm-anchors). mm-aliases.tex may name a finding, \mmalias{name}{25.91}; \mm{name} then prints and
% links the target section, so a finding that moves to a later section is re-pointed in one place.
\newcommand{\mmalias}[2]{\@ifundefined{mm@#2}{\PackageWarning{g17}{alias #1 points at undefined MM section #2}}{%
  \expandafter\edef\csname mm@#1\endcsname{\csname mm@#2\endcsname}%
  \expandafter\def\csname mmshow@#1\endcsname{#2}}}
\DeclareRobustCommand{\mm}[1]{%
  \@ifundefined{mm@#1}{\textbf{??#1}\PackageWarning{g17}{undefined MM section #1}}{%
  \href{\mmbase\#\csname mm@#1\endcsname}{\@ifundefined{mmshow@#1}{\mm@show#1/\@nil}%
    {\expandafter\expandafter\expandafter\mm@show\csname mmshow@#1\endcsname/\@nil}}}}
\makeatother
\input{mm-anchors}
% Named machine-model findings: \mmalias{name}{section}. \mm{name} prints and links the target section, so when a
% later section supersedes a finding, re-pointing the alias moves every citation of it. Section numbers remain the
% normal citation (\mm{25.91}); add an alias only for a finding likely to move. tools/g17docs.py --check refuses an
% alias whose target no machine-model heading defines, or that collides with a section number.


% A labelled box for an open question that qualifies the rule beside it.
\newtcolorbox{openq}[1][]{enhanced, breakable, colback=white, colframe=rule,
  boxrule=0.5pt, leftrule=2.5pt, colbacktitle=white, coltitle=native, arc=0pt,
  fonttitle=\small\bfseries, title={Open}, attach title to upper={\enspace}, #1}

% --- Figure styles -----------------------------------------------------------
\tikzset{
  >={Stealth[length=5pt]},
  box/.style={draw, rounded corners=2pt, align=center, font=\small, inner sep=5pt, line width=0.6pt,
    execute at begin node={\hyphenpenalty=10000\exhyphenpenalty=10000}},
  ours/.style={box, draw=ours, fill=ours!8, text=ink},
  nat/.style={box, draw=native, fill=native!8, text=ink},
  apple/.style={box, draw=apple, fill=apple!12, text=ink},
  band/.style={box, draw=apple, fill=apple!12, text=ink, minimum height=7mm},
  flow/.style={->, line width=0.7pt, color=ink!80},
  tag/.style={font=\scriptsize\itshape, text=ink!75, align=center},
}
